% Options for packages loaded elsewhere
\PassOptionsToPackage{unicode}{hyperref}
\PassOptionsToPackage{hyphens}{url}
\documentclass[
]{article}
\usepackage{xcolor}
\usepackage[margin=1in]{geometry}
\usepackage{amsmath,amssymb}
\setcounter{secnumdepth}{5}
\usepackage{iftex}
\ifPDFTeX
  \usepackage[T1]{fontenc}
  \usepackage[utf8]{inputenc}
  \usepackage{textcomp} % provide euro and other symbols
\else % if luatex or xetex
  \usepackage{unicode-math} % this also loads fontspec
  \defaultfontfeatures{Scale=MatchLowercase}
  \defaultfontfeatures[\rmfamily]{Ligatures=TeX,Scale=1}
\fi
\usepackage{lmodern}
\ifPDFTeX\else
  % xetex/luatex font selection
\fi
% Use upquote if available, for straight quotes in verbatim environments
\IfFileExists{upquote.sty}{\usepackage{upquote}}{}
\IfFileExists{microtype.sty}{% use microtype if available
  \usepackage[]{microtype}
  \UseMicrotypeSet[protrusion]{basicmath} % disable protrusion for tt fonts
}{}
\makeatletter
\@ifundefined{KOMAClassName}{% if non-KOMA class
  \IfFileExists{parskip.sty}{%
    \usepackage{parskip}
  }{% else
    \setlength{\parindent}{0pt}
    \setlength{\parskip}{6pt plus 2pt minus 1pt}}
}{% if KOMA class
  \KOMAoptions{parskip=half}}
\makeatother
\usepackage{color}
\usepackage{fancyvrb}
\newcommand{\VerbBar}{|}
\newcommand{\VERB}{\Verb[commandchars=\\\{\}]}
\DefineVerbatimEnvironment{Highlighting}{Verbatim}{commandchars=\\\{\}}
% Add ',fontsize=\small' for more characters per line
\usepackage{framed}
\definecolor{shadecolor}{RGB}{248,248,248}
\newenvironment{Shaded}{\begin{snugshade}}{\end{snugshade}}
\newcommand{\AlertTok}[1]{\textcolor[rgb]{0.94,0.16,0.16}{#1}}
\newcommand{\AnnotationTok}[1]{\textcolor[rgb]{0.56,0.35,0.01}{\textbf{\textit{#1}}}}
\newcommand{\AttributeTok}[1]{\textcolor[rgb]{0.13,0.29,0.53}{#1}}
\newcommand{\BaseNTok}[1]{\textcolor[rgb]{0.00,0.00,0.81}{#1}}
\newcommand{\BuiltInTok}[1]{#1}
\newcommand{\CharTok}[1]{\textcolor[rgb]{0.31,0.60,0.02}{#1}}
\newcommand{\CommentTok}[1]{\textcolor[rgb]{0.56,0.35,0.01}{\textit{#1}}}
\newcommand{\CommentVarTok}[1]{\textcolor[rgb]{0.56,0.35,0.01}{\textbf{\textit{#1}}}}
\newcommand{\ConstantTok}[1]{\textcolor[rgb]{0.56,0.35,0.01}{#1}}
\newcommand{\ControlFlowTok}[1]{\textcolor[rgb]{0.13,0.29,0.53}{\textbf{#1}}}
\newcommand{\DataTypeTok}[1]{\textcolor[rgb]{0.13,0.29,0.53}{#1}}
\newcommand{\DecValTok}[1]{\textcolor[rgb]{0.00,0.00,0.81}{#1}}
\newcommand{\DocumentationTok}[1]{\textcolor[rgb]{0.56,0.35,0.01}{\textbf{\textit{#1}}}}
\newcommand{\ErrorTok}[1]{\textcolor[rgb]{0.64,0.00,0.00}{\textbf{#1}}}
\newcommand{\ExtensionTok}[1]{#1}
\newcommand{\FloatTok}[1]{\textcolor[rgb]{0.00,0.00,0.81}{#1}}
\newcommand{\FunctionTok}[1]{\textcolor[rgb]{0.13,0.29,0.53}{\textbf{#1}}}
\newcommand{\ImportTok}[1]{#1}
\newcommand{\InformationTok}[1]{\textcolor[rgb]{0.56,0.35,0.01}{\textbf{\textit{#1}}}}
\newcommand{\KeywordTok}[1]{\textcolor[rgb]{0.13,0.29,0.53}{\textbf{#1}}}
\newcommand{\NormalTok}[1]{#1}
\newcommand{\OperatorTok}[1]{\textcolor[rgb]{0.81,0.36,0.00}{\textbf{#1}}}
\newcommand{\OtherTok}[1]{\textcolor[rgb]{0.56,0.35,0.01}{#1}}
\newcommand{\PreprocessorTok}[1]{\textcolor[rgb]{0.56,0.35,0.01}{\textit{#1}}}
\newcommand{\RegionMarkerTok}[1]{#1}
\newcommand{\SpecialCharTok}[1]{\textcolor[rgb]{0.81,0.36,0.00}{\textbf{#1}}}
\newcommand{\SpecialStringTok}[1]{\textcolor[rgb]{0.31,0.60,0.02}{#1}}
\newcommand{\StringTok}[1]{\textcolor[rgb]{0.31,0.60,0.02}{#1}}
\newcommand{\VariableTok}[1]{\textcolor[rgb]{0.00,0.00,0.00}{#1}}
\newcommand{\VerbatimStringTok}[1]{\textcolor[rgb]{0.31,0.60,0.02}{#1}}
\newcommand{\WarningTok}[1]{\textcolor[rgb]{0.56,0.35,0.01}{\textbf{\textit{#1}}}}
\usepackage{longtable,booktabs,array}
\usepackage{calc} % for calculating minipage widths
% Correct order of tables after \paragraph or \subparagraph
\usepackage{etoolbox}
\makeatletter
\patchcmd\longtable{\par}{\if@noskipsec\mbox{}\fi\par}{}{}
\makeatother
% Allow footnotes in longtable head/foot
\IfFileExists{footnotehyper.sty}{\usepackage{footnotehyper}}{\usepackage{footnote}}
\makesavenoteenv{longtable}
\usepackage{graphicx}
\makeatletter
\newsavebox\pandoc@box
\newcommand*\pandocbounded[1]{% scales image to fit in text height/width
  \sbox\pandoc@box{#1}%
  \Gscale@div\@tempa{\textheight}{\dimexpr\ht\pandoc@box+\dp\pandoc@box\relax}%
  \Gscale@div\@tempb{\linewidth}{\wd\pandoc@box}%
  \ifdim\@tempb\p@<\@tempa\p@\let\@tempa\@tempb\fi% select the smaller of both
  \ifdim\@tempa\p@<\p@\scalebox{\@tempa}{\usebox\pandoc@box}%
  \else\usebox{\pandoc@box}%
  \fi%
}
% Set default figure placement to htbp
\def\fps@figure{htbp}
\makeatother
\setlength{\emergencystretch}{3em} % prevent overfull lines
\providecommand{\tightlist}{%
  \setlength{\itemsep}{0pt}\setlength{\parskip}{0pt}}
\usepackage{bookmark}
\IfFileExists{xurl.sty}{\usepackage{xurl}}{} % add URL line breaks if available
\urlstyle{same}
\hypersetup{
  pdftitle={Inteligência Artificial Local com Ollama},
  pdfauthor={Paulo Pimenta},
  hidelinks,
  pdfcreator={LaTeX via pandoc}}

\title{Inteligência Artificial Local com Ollama}
\usepackage{etoolbox}
\makeatletter
\providecommand{\subtitle}[1]{% add subtitle to \maketitle
  \apptocmd{\@title}{\par {\large #1 \par}}{}{}
}
\makeatother
\subtitle{Um guia prático do zero ao servidor de IA}
\author{Paulo Pimenta}
\date{2026-05-02}

\begin{document}
\maketitle

{
\setcounter{tocdepth}{3}
\tableofcontents
}
\begin{center}\rule{0.5\linewidth}{0.5pt}\end{center}

\section{Introdução}\label{introduuxe7uxe3o}

Este tutorial apresenta de forma prática e objetiva como utilizar
modelos de inteligência artificial (IA) localmente, sem depender de
serviços em nuvem. Abordaremos desde os conceitos fundamentais até a
criação de um servidor de IA acessível por múltiplos dispositivos na
mesma rede.

A execução local de modelos de linguagem (LLMs) oferece vantagens
importantes:

\begin{itemize}
\tightlist
\item
  \textbf{Privacidade}: os dados não saem do seu computador
\item
  \textbf{Custo zero}: sem limites de uso ou assinaturas
\item
  \textbf{Disponibilidade offline}: funciona sem internet
\item
  \textbf{Controle total}: você escolhe e gerencia os modelos
\end{itemize}

\begin{center}\rule{0.5\linewidth}{0.5pt}\end{center}

\section{O Ollama e os Modelos de IA}\label{o-ollama-e-os-modelos-de-ia}

\subsection{O que é o Ollama?}\label{o-que-uxe9-o-ollama}

O \textbf{Ollama} é uma plataforma de código aberto que permite baixar,
gerenciar e executar modelos de linguagem grandes (LLMs) localmente,
diretamente no seu computador. Ele abstrai toda a complexidade técnica
de executar esses modelos, oferecendo uma interface simples via linha de
comando e uma API REST integrada.

Em termos práticos, o Ollama funciona como um gerenciador de modelos de
IA --- de modo similar ao que o \texttt{pip} é para pacotes Python ou o
\texttt{apt} é para pacotes no Ubuntu.

\subsection{O que são Modelos de Linguagem
(LLMs)?}\label{o-que-suxe3o-modelos-de-linguagem-llms}

Modelos de linguagem grandes (do inglês \emph{Large Language Models},
LLMs) são sistemas de IA treinados em enormes volumes de texto com o
objetivo de compreender e gerar linguagem natural. Eles são capazes de:

\begin{itemize}
\tightlist
\item
  Responder perguntas e manter conversas
\item
  Escrever, revisar e explicar código
\item
  Resumir e analisar documentos
\item
  Auxiliar em tarefas de ciência de dados
\end{itemize}

\subsection{Principais modelos disponíveis no
Ollama}\label{principais-modelos-disponuxedveis-no-ollama}

A tabela a seguir apresenta os modelos mais relevantes disponíveis via
Ollama, organizados por finalidade e tamanho:

\begin{longtable}[]{@{}
  >{\raggedright\arraybackslash}p{(\linewidth - 6\tabcolsep) * \real{0.2500}}
  >{\raggedright\arraybackslash}p{(\linewidth - 6\tabcolsep) * \real{0.2500}}
  >{\raggedright\arraybackslash}p{(\linewidth - 6\tabcolsep) * \real{0.2500}}
  >{\raggedright\arraybackslash}p{(\linewidth - 6\tabcolsep) * \real{0.2500}}@{}}
\toprule\noalign{}
\begin{minipage}[b]{\linewidth}\raggedright
Modelo
\end{minipage} & \begin{minipage}[b]{\linewidth}\raggedright
Tamanho
\end{minipage} & \begin{minipage}[b]{\linewidth}\raggedright
Especialidade
\end{minipage} & \begin{minipage}[b]{\linewidth}\raggedright
Indicado para
\end{minipage} \\
\midrule\noalign{}
\endhead
\bottomrule\noalign{}
\endlastfoot
\texttt{tinyllama} & 637 MB & Uso geral leve & Hardware limitado \\
\texttt{qwen2.5-coder:0.5b} & 400 MB & Código & Hardware limitado \\
\texttt{qwen2.5-coder:1.5b} & 1.0 GB & Código e dados & Hardware
intermediário \\
\texttt{qwen2.5-coder:7b} & 4.7 GB & Código avançado & Hardware
potente \\
\texttt{llama3.2:1b} & 1.3 GB & Uso geral & Hardware intermediário \\
\texttt{llama3.2:latest} & 2.0 GB & Uso geral & Hardware
intermediário \\
\texttt{phi3:mini} & 2.3 GB & Raciocínio & Hardware intermediário \\
\texttt{deepseek-coder:6.7b} & 3.8 GB & Código avançado & Hardware
potente \\
\texttt{mistral:7b} & 4.1 GB & Uso geral avançado & Hardware potente \\
\end{longtable}

\begin{quote}
\textbf{Nota}: o tamanho do modelo influencia diretamente a qualidade
das respostas e os recursos de hardware necessários para executá-lo.
Modelos maiores tendem a ser mais capazes, porém mais lentos em hardware
limitado.
\end{quote}

\subsection{Tipos de modelos por
finalidade}\label{tipos-de-modelos-por-finalidade}

\subsubsection{Modelos generalistas}\label{modelos-generalistas}

São treinados para tarefas diversas como conversação, redação, resumo e
raciocínio geral. Exemplos: \texttt{llama3.2}, \texttt{tinyllama},
\texttt{mistral}.

\subsubsection{Modelos especialistas em
código}\label{modelos-especialistas-em-cuxf3digo}

Treinados com grande volume de código-fonte em múltiplas linguagens. São
superiores aos generalistas para escrever, depurar e explicar código.
Exemplos: \texttt{qwen2.5-coder}, \texttt{deepseek-coder},
\texttt{codellama}.

\begin{center}\rule{0.5\linewidth}{0.5pt}\end{center}

\section{Hardware e Configuração dos
Modelos}\label{hardware-e-configurauxe7uxe3o-dos-modelos}

\subsection{Por que o hardware
importa?}\label{por-que-o-hardware-importa}

A execução local de LLMs é intensiva em recursos computacionais. O
hardware determina diretamente:

\begin{itemize}
\tightlist
\item
  \textbf{Quais modelos} você consegue executar
\item
  \textbf{A velocidade} de geração de respostas (tokens por segundo)
\item
  \textbf{A qualidade} máxima que você pode alcançar
\end{itemize}

\subsection{Componentes críticos}\label{componentes-cruxedticos}

\subsubsection{GPU (Placa de Vídeo)}\label{gpu-placa-de-vuxeddeo}

É o componente mais importante para executar LLMs com eficiência. GPUs
possuem centenas ou milhares de núcleos paralelos, ideais para as
operações matriciais exigidas pelos modelos.

\begin{itemize}
\tightlist
\item
  \textbf{NVIDIA}: melhor suporte via CUDA --- recomendada para uso com
  Ollama
\item
  \textbf{AMD}: suporte via ROCm (mais limitado no Windows)
\item
  \textbf{GPU integrada}: não é aproveitada pelo Ollama na maioria dos
  casos
\end{itemize}

\subsubsection{RAM}\label{ram}

Quando o modelo não cabe inteiramente na VRAM da GPU, o excedente é
carregado na RAM do sistema. Quanto mais RAM disponível, maiores os
modelos que podem ser executados.

\subsubsection{CPU}\label{cpu}

Utilizada quando não há GPU compatível. CPUs modernas executam LLMs de
forma funcional, porém significativamente mais lenta que GPUs dedicadas.

\subsection{Tabela de recomendações por
hardware}\label{tabela-de-recomendauxe7uxf5es-por-hardware}

\begin{longtable}[]{@{}
  >{\raggedright\arraybackslash}p{(\linewidth - 6\tabcolsep) * \real{0.2500}}
  >{\raggedright\arraybackslash}p{(\linewidth - 6\tabcolsep) * \real{0.2500}}
  >{\raggedright\arraybackslash}p{(\linewidth - 6\tabcolsep) * \real{0.2500}}
  >{\raggedright\arraybackslash}p{(\linewidth - 6\tabcolsep) * \real{0.2500}}@{}}
\toprule\noalign{}
\begin{minipage}[b]{\linewidth}\raggedright
Configuração
\end{minipage} & \begin{minipage}[b]{\linewidth}\raggedright
RAM
\end{minipage} & \begin{minipage}[b]{\linewidth}\raggedright
VRAM
\end{minipage} & \begin{minipage}[b]{\linewidth}\raggedright
Modelos recomendados
\end{minipage} \\
\midrule\noalign{}
\endhead
\bottomrule\noalign{}
\endlastfoot
Hardware muito limitado & 8 GB & Sem GPU dedicada & \texttt{tinyllama},
\texttt{qwen2.5-coder:0.5b} \\
Hardware intermediário & 16 GB & 2--4 GB & \texttt{qwen2.5-coder:1.5b},
\texttt{llama3.2:1b}, \texttt{phi3:mini} \\
Hardware bom & 16 GB & 6--8 GB & \texttt{qwen2.5-coder:7b},
\texttt{deepseek-coder:6.7b} \\
Hardware avançado & 32 GB+ & 12 GB+ & \texttt{llama3.1:70b},
\texttt{deepseek-coder-v2} \\
\end{longtable}

\subsection{Exemplos práticos de
configuração}\label{exemplos-pruxe1ticos-de-configurauxe7uxe3o}

\subsubsection{Notebook com CPU antiga e GPU
integrada}\label{notebook-com-cpu-antiga-e-gpu-integrada}

\textbf{Exemplo}: HP Pavilion com AMD A10-5745M (2013), 16 GB RAM, GPU
Radeon HD 8610G integrada.

\begin{itemize}
\tightlist
\item
  Sem suporte a CUDA ou ROCm moderno
\item
  Execução apenas na CPU --- velocidade limitada (\textasciitilde2--5
  tokens/segundo)
\item
  Modelos recomendados: \texttt{qwen2.5-coder:0.5b} e \texttt{tinyllama}
\end{itemize}

\subsubsection{Notebook moderno com GPU dedicada de
entrada}\label{notebook-moderno-com-gpu-dedicada-de-entrada}

\textbf{Exemplo}: Dell com Intel Core i7-1165G7, 16 GB RAM, NVIDIA MX350
(2 GB VRAM).

\begin{itemize}
\tightlist
\item
  Suporte a CUDA --- Ollama usa a GPU automaticamente
\item
  Execução híbrida GPU + RAM para modelos até \textasciitilde5 GB
\item
  Velocidade estimada: 10--20 tokens/segundo
\item
  Modelos recomendados: \texttt{qwen2.5-coder:1.5b},
  \texttt{llama3.2:latest}, \texttt{phi3:mini}
\end{itemize}

\begin{center}\rule{0.5\linewidth}{0.5pt}\end{center}

\section{Instalando o Ollama}\label{instalando-o-ollama}

\subsection{Instalação no Windows}\label{instalauxe7uxe3o-no-windows}

\subsubsection{Requisitos}\label{requisitos}

\begin{itemize}
\tightlist
\item
  Windows 10 ou superior (64 bits)
\item
  Pelo menos 8 GB de RAM
\item
  Espaço em disco suficiente para os modelos desejados
\end{itemize}

\subsubsection{Passo a passo}\label{passo-a-passo}

\textbf{1. Baixar o instalador}

Acesse \url{https://ollama.com/download} e clique em \textbf{Download
for Windows}.

\textbf{2. Executar o instalador}

Execute o arquivo \texttt{.exe} baixado e siga as instruções da tela. O
Ollama será instalado e iniciado automaticamente como um serviço em
segundo plano.

\textbf{3. Verificar a instalação}

Abra o \textbf{PowerShell} e execute:

\begin{Shaded}
\begin{Highlighting}[]
\ExtensionTok{ollama} \AttributeTok{{-}{-}version}
\end{Highlighting}
\end{Shaded}

Se a instalação foi bem-sucedida, você verá a versão instalada, por
exemplo:

\begin{verbatim}
ollama version 0.3.12
\end{verbatim}

\textbf{4. Verificar se o serviço está ativo}

\begin{Shaded}
\begin{Highlighting}[]
\CommentTok{\# No navegador ou PowerShell}
\ExtensionTok{curl}\NormalTok{ http://localhost:11434}
\end{Highlighting}
\end{Shaded}

A resposta esperada é:

\begin{verbatim}
Ollama is running
\end{verbatim}

\subsection{Instalação no Linux (Ubuntu 22.04
LTS)}\label{instalauxe7uxe3o-no-linux-ubuntu-22.04-lts}

\subsubsection{Requisitos}\label{requisitos-1}

\begin{itemize}
\tightlist
\item
  Ubuntu 20.04 ou superior
\item
  Pelo menos 8 GB de RAM
\item
  Curl instalado
\end{itemize}

\subsubsection{Passo a passo}\label{passo-a-passo-1}

\textbf{1. Instalar via script oficial}

\begin{Shaded}
\begin{Highlighting}[]
\ExtensionTok{curl} \AttributeTok{{-}fsSL}\NormalTok{ https://ollama.com/install.sh }\KeywordTok{|} \FunctionTok{sh}
\end{Highlighting}
\end{Shaded}

O script detecta automaticamente a arquitetura do sistema e instala a
versão adequada.

\textbf{2. Verificar a instalação}

\begin{Shaded}
\begin{Highlighting}[]
\ExtensionTok{ollama} \AttributeTok{{-}{-}version}
\end{Highlighting}
\end{Shaded}

\textbf{3. Verificar se o serviço está ativo}

\begin{Shaded}
\begin{Highlighting}[]
\ExtensionTok{systemctl}\NormalTok{ status ollama}
\end{Highlighting}
\end{Shaded}

A saída deve mostrar \texttt{active\ (running)}.

\textbf{4. Iniciar o serviço manualmente (se necessário)}

\begin{Shaded}
\begin{Highlighting}[]
\ExtensionTok{ollama}\NormalTok{ serve}
\end{Highlighting}
\end{Shaded}

\begin{quote}
\textbf{Nota}: em instalações via script oficial, o Ollama é registrado
como serviço do sistema e iniciado automaticamente. Caso tenha instalado
via Snap ou Flatpak, o diretório dos modelos pode estar em um caminho
diferente do padrão (\texttt{\textasciitilde{}/.ollama/models}).
\end{quote}

\begin{center}\rule{0.5\linewidth}{0.5pt}\end{center}

\section{Instalando Modelos pelo
Ollama}\label{instalando-modelos-pelo-ollama}

\subsection{\texorpdfstring{O comando
\texttt{ollama\ pull}}{O comando ollama pull}}\label{o-comando-ollama-pull}

Para baixar um modelo, utilize o comando \texttt{ollama\ pull} seguido
do nome do modelo:

\begin{Shaded}
\begin{Highlighting}[]
\ExtensionTok{ollama}\NormalTok{ pull }\OperatorTok{\textless{}}\NormalTok{nome\_do\_modelo}\OperatorTok{\textgreater{}}
\end{Highlighting}
\end{Shaded}

\subsection{Exemplos práticos}\label{exemplos-pruxe1ticos}

\subsubsection{Windows (PowerShell)}\label{windows-powershell}

\begin{Shaded}
\begin{Highlighting}[]
\CommentTok{\# Baixar o TinyLlama (leve, ideal para hardware limitado)}
\ExtensionTok{ollama}\NormalTok{ pull tinyllama}

\CommentTok{\# Baixar o Qwen 2.5 Coder 0.5b (especialista em código, muito leve)}
\ExtensionTok{ollama}\NormalTok{ pull qwen2.5{-}coder:0.5b}

\CommentTok{\# Baixar o Qwen 2.5 Coder 1.5b (especialista em código, intermediário)}
\ExtensionTok{ollama}\NormalTok{ pull qwen2.5{-}coder:1.5b}
\end{Highlighting}
\end{Shaded}

\subsubsection{Linux (Terminal)}\label{linux-terminal}

\begin{Shaded}
\begin{Highlighting}[]
\CommentTok{\# Os comandos são idênticos ao Windows}
\ExtensionTok{ollama}\NormalTok{ pull tinyllama}
\ExtensionTok{ollama}\NormalTok{ pull qwen2.5{-}coder:0.5b}
\end{Highlighting}
\end{Shaded}

\subsection{Gerenciando modelos
instalados}\label{gerenciando-modelos-instalados}

\begin{Shaded}
\begin{Highlighting}[]
\CommentTok{\# Listar todos os modelos instalados}
\ExtensionTok{ollama}\NormalTok{ list}

\CommentTok{\# Remover um modelo (libera espaço em disco)}
\ExtensionTok{ollama}\NormalTok{ rm tinyllama}

\CommentTok{\# Exibir informações detalhadas de um modelo}
\ExtensionTok{ollama}\NormalTok{ show qwen2.5{-}coder:0.5b}
\end{Highlighting}
\end{Shaded}

\begin{quote}
\textbf{Importante}: o comando \texttt{ollama\ rm} remove completamente
os arquivos do modelo do disco. Não é necessário remover diretórios
manualmente --- o Ollama gerencia tudo internamente.
\end{quote}

\subsection{Monitorando o espaço em
disco}\label{monitorando-o-espauxe7o-em-disco}

\subsubsection{Windows}\label{windows}

\begin{Shaded}
\begin{Highlighting}[]
\CommentTok{\# Ver espaço disponível em disco}
\ExtensionTok{Get{-}PSDrive} \AttributeTok{{-}PSProvider}\NormalTok{ FileSystem}
\end{Highlighting}
\end{Shaded}

\subsubsection{Linux}\label{linux}

\begin{Shaded}
\begin{Highlighting}[]
\CommentTok{\# Ver espaço disponível em disco}
\FunctionTok{df} \AttributeTok{{-}h}
\end{Highlighting}
\end{Shaded}

\begin{center}\rule{0.5\linewidth}{0.5pt}\end{center}

\section{Usando os Modelos}\label{usando-os-modelos}

\subsection{Uso via linha de comando}\label{uso-via-linha-de-comando}

A forma mais direta de interagir com um modelo é pelo terminal:

\begin{Shaded}
\begin{Highlighting}[]
\CommentTok{\# Iniciar uma conversa interativa}
\ExtensionTok{ollama}\NormalTok{ run tinyllama}

\CommentTok{\# Enviar uma pergunta direta (sem modo interativo)}
\ExtensionTok{ollama}\NormalTok{ run qwen2.5{-}coder:0.5b }\StringTok{"Como fazer um loop em Python?"}
\end{Highlighting}
\end{Shaded}

Para encerrar o modo interativo, digite \texttt{/bye} ou pressione
\texttt{Ctrl+D}.

\subsection{Uso via Python}\label{uso-via-python}

\subsubsection{Instalação da
biblioteca}\label{instalauxe7uxe3o-da-biblioteca}

\begin{Shaded}
\begin{Highlighting}[]
\ExtensionTok{pip}\NormalTok{ install ollama}
\end{Highlighting}
\end{Shaded}

\subsubsection{Exemplo básico}\label{exemplo-buxe1sico}

\begin{Shaded}
\begin{Highlighting}[]
\ImportTok{import}\NormalTok{ ollama}

\NormalTok{response }\OperatorTok{=}\NormalTok{ ollama.chat(}
\NormalTok{    model}\OperatorTok{=}\StringTok{\textquotesingle{}qwen2.5{-}coder:0.5b\textquotesingle{}}\NormalTok{,}
\NormalTok{    messages}\OperatorTok{=}\NormalTok{[}
\NormalTok{        \{}\StringTok{\textquotesingle{}role\textquotesingle{}}\NormalTok{: }\StringTok{\textquotesingle{}user\textquotesingle{}}\NormalTok{, }\StringTok{\textquotesingle{}content\textquotesingle{}}\NormalTok{: }\StringTok{\textquotesingle{}Explique o que é regressão linear.\textquotesingle{}}\NormalTok{\}}
\NormalTok{    ]}
\NormalTok{)}

\BuiltInTok{print}\NormalTok{(response[}\StringTok{\textquotesingle{}message\textquotesingle{}}\NormalTok{][}\StringTok{\textquotesingle{}content\textquotesingle{}}\NormalTok{])}
\end{Highlighting}
\end{Shaded}

\subsubsection{Chat com histórico de
conversa}\label{chat-com-histuxf3rico-de-conversa}

\begin{Shaded}
\begin{Highlighting}[]
\ImportTok{import}\NormalTok{ ollama}

\NormalTok{historico }\OperatorTok{=}\NormalTok{ []}

\BuiltInTok{print}\NormalTok{(}\StringTok{"Chat iniciado! Digite \textquotesingle{}sair\textquotesingle{} para encerrar.}\CharTok{\textbackslash{}n}\StringTok{"}\NormalTok{)}

\ControlFlowTok{while} \VariableTok{True}\NormalTok{:}
\NormalTok{    pergunta }\OperatorTok{=} \BuiltInTok{input}\NormalTok{(}\StringTok{"Você: "}\NormalTok{)}

    \ControlFlowTok{if}\NormalTok{ pergunta.lower() }\OperatorTok{==} \StringTok{\textquotesingle{}sair\textquotesingle{}}\NormalTok{:}
        \BuiltInTok{print}\NormalTok{(}\StringTok{"Encerrando..."}\NormalTok{)}
        \ControlFlowTok{break}

\NormalTok{    historico.append(\{}\StringTok{\textquotesingle{}role\textquotesingle{}}\NormalTok{: }\StringTok{\textquotesingle{}user\textquotesingle{}}\NormalTok{, }\StringTok{\textquotesingle{}content\textquotesingle{}}\NormalTok{: pergunta\})}

\NormalTok{    response }\OperatorTok{=}\NormalTok{ ollama.chat(}
\NormalTok{        model}\OperatorTok{=}\StringTok{\textquotesingle{}qwen2.5{-}coder:0.5b\textquotesingle{}}\NormalTok{,}
\NormalTok{        messages}\OperatorTok{=}\NormalTok{historico}
\NormalTok{    )}

\NormalTok{    resposta }\OperatorTok{=}\NormalTok{ response[}\StringTok{\textquotesingle{}message\textquotesingle{}}\NormalTok{][}\StringTok{\textquotesingle{}content\textquotesingle{}}\NormalTok{]}
\NormalTok{    historico.append(\{}\StringTok{\textquotesingle{}role\textquotesingle{}}\NormalTok{: }\StringTok{\textquotesingle{}assistant\textquotesingle{}}\NormalTok{, }\StringTok{\textquotesingle{}content\textquotesingle{}}\NormalTok{: resposta\})}

    \BuiltInTok{print}\NormalTok{(}\SpecialStringTok{f"}\CharTok{\textbackslash{}n}\SpecialStringTok{Modelo: }\SpecialCharTok{\{}\NormalTok{resposta}\SpecialCharTok{\}}\CharTok{\textbackslash{}n}\SpecialStringTok{"}\NormalTok{)}
\end{Highlighting}
\end{Shaded}

\subsubsection{Streaming de respostas}\label{streaming-de-respostas}

O streaming exibe a resposta progressivamente, melhorando a experiência
em hardware mais lento:

\begin{Shaded}
\begin{Highlighting}[]
\ImportTok{import}\NormalTok{ ollama}

\ControlFlowTok{for}\NormalTok{ chunk }\KeywordTok{in}\NormalTok{ ollama.chat(}
\NormalTok{    model}\OperatorTok{=}\StringTok{\textquotesingle{}qwen2.5{-}coder:0.5b\textquotesingle{}}\NormalTok{,}
\NormalTok{    messages}\OperatorTok{=}\NormalTok{[\{}\StringTok{\textquotesingle{}role\textquotesingle{}}\NormalTok{: }\StringTok{\textquotesingle{}user\textquotesingle{}}\NormalTok{, }\StringTok{\textquotesingle{}content\textquotesingle{}}\NormalTok{: }\StringTok{\textquotesingle{}O que é machine learning?\textquotesingle{}}\NormalTok{\}],}
\NormalTok{    stream}\OperatorTok{=}\VariableTok{True}
\NormalTok{):}
    \BuiltInTok{print}\NormalTok{(chunk[}\StringTok{\textquotesingle{}message\textquotesingle{}}\NormalTok{][}\StringTok{\textquotesingle{}content\textquotesingle{}}\NormalTok{], end}\OperatorTok{=}\StringTok{\textquotesingle{}\textquotesingle{}}\NormalTok{, flush}\OperatorTok{=}\VariableTok{True}\NormalTok{)}
\end{Highlighting}
\end{Shaded}

\subsubsection{Exemplo aplicado à ciência de
dados}\label{exemplo-aplicado-uxe0-ciuxeancia-de-dados}

\begin{Shaded}
\begin{Highlighting}[]
\ImportTok{import}\NormalTok{ ollama}

\NormalTok{pergunta }\OperatorTok{=} \StringTok{"""}
\StringTok{Em Python, usando pandas e matplotlib, como faço para:}
\StringTok{1. Carregar um CSV}
\StringTok{2. Calcular estatísticas descritivas}
\StringTok{3. Plotar um histograma da coluna \textquotesingle{}idade\textquotesingle{}}
\StringTok{"""}

\ControlFlowTok{for}\NormalTok{ chunk }\KeywordTok{in}\NormalTok{ ollama.chat(}
\NormalTok{    model}\OperatorTok{=}\StringTok{\textquotesingle{}qwen2.5{-}coder:0.5b\textquotesingle{}}\NormalTok{,}
\NormalTok{    messages}\OperatorTok{=}\NormalTok{[\{}\StringTok{\textquotesingle{}role\textquotesingle{}}\NormalTok{: }\StringTok{\textquotesingle{}user\textquotesingle{}}\NormalTok{, }\StringTok{\textquotesingle{}content\textquotesingle{}}\NormalTok{: pergunta\}],}
\NormalTok{    stream}\OperatorTok{=}\VariableTok{True}
\NormalTok{):}
    \BuiltInTok{print}\NormalTok{(chunk[}\StringTok{\textquotesingle{}message\textquotesingle{}}\NormalTok{][}\StringTok{\textquotesingle{}content\textquotesingle{}}\NormalTok{], end}\OperatorTok{=}\StringTok{\textquotesingle{}\textquotesingle{}}\NormalTok{, flush}\OperatorTok{=}\VariableTok{True}\NormalTok{)}
\end{Highlighting}
\end{Shaded}

\subsection{Uso via API REST}\label{uso-via-api-rest}

O Ollama expõe uma API REST na porta 11434, compatível com o formato da
API da OpenAI:

\begin{Shaded}
\begin{Highlighting}[]
\CommentTok{\# Enviar uma requisição via curl}
\ExtensionTok{curl}\NormalTok{ http://localhost:11434/api/chat }\DataTypeTok{\textbackslash{}}
  \AttributeTok{{-}H} \StringTok{"Content{-}Type: application/json"} \DataTypeTok{\textbackslash{}}
  \AttributeTok{{-}d} \StringTok{\textquotesingle{}\{}
\StringTok{    "model": "qwen2.5{-}coder:0.5b",}
\StringTok{    "messages": [}
\StringTok{      \{"role": "user", "content": "O que é uma função recursiva?"\}}
\StringTok{    ],}
\StringTok{    "stream": false}
\StringTok{  \}\textquotesingle{}}
\end{Highlighting}
\end{Shaded}

\begin{center}\rule{0.5\linewidth}{0.5pt}\end{center}

\section{Integrando Modelos nas IDEs com o Plugin
Continue}\label{integrando-modelos-nas-ides-com-o-plugin-continue}

\subsection{O que é o Continue?}\label{o-que-uxe9-o-continue}

O \textbf{Continue} é um plugin de código aberto para IDEs que integra
modelos de IA diretamente no ambiente de desenvolvimento. Ele suporta VS
Code e toda a família JetBrains (PyCharm, IntelliJ, etc.) e permite usar
modelos locais via Ollama sem custo e sem enviar código para servidores
externos.

\subsection{Instalação no VS Code}\label{instalauxe7uxe3o-no-vs-code}

\textbf{1.} Abra o VS Code e pressione \texttt{Ctrl+Shift+X}

\textbf{2.} Pesquise por \textbf{``Continue''} na barra de busca

\textbf{3.} Clique em \textbf{Install}

\textbf{4.} Após a instalação, o ícone do Continue aparecerá na barra
lateral esquerda

\subsection{\texorpdfstring{Configuração do arquivo
\texttt{config.yaml}}{Configuração do arquivo config.yaml}}\label{configurauxe7uxe3o-do-arquivo-config.yaml}

O Continue é configurado pelo arquivo \texttt{config.yaml}, localizado
em:

\begin{itemize}
\tightlist
\item
  \textbf{Windows}:
  \texttt{C:\textbackslash{}Users\textbackslash{}\textless{}seu\_usuario\textgreater{}\textbackslash{}.continue\textbackslash{}config.yaml}
\item
  \textbf{Linux/macOS}: \texttt{\textasciitilde{}/.continue/config.yaml}
\end{itemize}

Para abrir o arquivo diretamente pelo terminal do VS Code
(\texttt{Ctrl+}):

\begin{Shaded}
\begin{Highlighting}[]
\CommentTok{\# Windows}
\ExtensionTok{code} \VariableTok{$env}\NormalTok{:USERPROFILE}\DataTypeTok{\textbackslash{}.}\NormalTok{continue}\DataTypeTok{\textbackslash{}c}\NormalTok{onfig.yaml}

\CommentTok{\# Linux}
\ExtensionTok{code}\NormalTok{ \textasciitilde{}/.continue/config.yaml}
\end{Highlighting}
\end{Shaded}

\subsubsection{Configuração com um único
modelo}\label{configurauxe7uxe3o-com-um-uxfanico-modelo}

\begin{Shaded}
\begin{Highlighting}[]
\FunctionTok{name}\KeywordTok{:}\AttributeTok{ Local Config}
\FunctionTok{version}\KeywordTok{:}\AttributeTok{ }\FloatTok{1.0.0}
\FunctionTok{schema}\KeywordTok{:}\AttributeTok{ v1}

\FunctionTok{models}\KeywordTok{:}
\AttributeTok{  }\KeywordTok{{-}}\AttributeTok{ }\FunctionTok{name}\KeywordTok{:}\AttributeTok{ Qwen 2.5 Coder 0.5b}
\AttributeTok{    }\FunctionTok{provider}\KeywordTok{:}\AttributeTok{ ollama}
\AttributeTok{    }\FunctionTok{model}\KeywordTok{:}\AttributeTok{ qwen2.5{-}coder:0.5b}
\AttributeTok{    }\FunctionTok{apiBase}\KeywordTok{:}\AttributeTok{ http://localhost:11434}

\FunctionTok{allowAnonymousTelemetry}\KeywordTok{:}\AttributeTok{ }\CharTok{false}
\end{Highlighting}
\end{Shaded}

\subsubsection{Configuração com múltiplos
modelos}\label{configurauxe7uxe3o-com-muxfaltiplos-modelos}

\begin{Shaded}
\begin{Highlighting}[]
\FunctionTok{name}\KeywordTok{:}\AttributeTok{ Local Config}
\FunctionTok{version}\KeywordTok{:}\AttributeTok{ }\FloatTok{1.0.0}
\FunctionTok{schema}\KeywordTok{:}\AttributeTok{ v1}

\FunctionTok{models}\KeywordTok{:}
\AttributeTok{  }\KeywordTok{{-}}\AttributeTok{ }\FunctionTok{name}\KeywordTok{:}\AttributeTok{ Qwen 2.5 Coder 0.5b}
\AttributeTok{    }\FunctionTok{provider}\KeywordTok{:}\AttributeTok{ ollama}
\AttributeTok{    }\FunctionTok{model}\KeywordTok{:}\AttributeTok{ qwen2.5{-}coder:0.5b}
\AttributeTok{    }\FunctionTok{apiBase}\KeywordTok{:}\AttributeTok{ http://localhost:11434}

\AttributeTok{  }\KeywordTok{{-}}\AttributeTok{ }\FunctionTok{name}\KeywordTok{:}\AttributeTok{ TinyLlama}
\AttributeTok{    }\FunctionTok{provider}\KeywordTok{:}\AttributeTok{ ollama}
\AttributeTok{    }\FunctionTok{model}\KeywordTok{:}\AttributeTok{ tinyllama:latest}
\AttributeTok{    }\FunctionTok{apiBase}\KeywordTok{:}\AttributeTok{ http://localhost:11434}

\FunctionTok{allowAnonymousTelemetry}\KeywordTok{:}\AttributeTok{ }\CharTok{false}
\end{Highlighting}
\end{Shaded}

\begin{quote}
\textbf{Atenção}: o campo \texttt{tabAutocompleteModel} pode não ser
suportado em versões mais recentes do Continue (1.2+). Nesse caso,
configure o autocomplete diretamente pelo painel da extensão.
\end{quote}

\subsection{Explicação dos campos do
config.yaml}\label{explicauxe7uxe3o-dos-campos-do-config.yaml}

\begin{longtable}[]{@{}
  >{\raggedright\arraybackslash}p{(\linewidth - 2\tabcolsep) * \real{0.5000}}
  >{\raggedright\arraybackslash}p{(\linewidth - 2\tabcolsep) * \real{0.5000}}@{}}
\toprule\noalign{}
\begin{minipage}[b]{\linewidth}\raggedright
Campo
\end{minipage} & \begin{minipage}[b]{\linewidth}\raggedright
Descrição
\end{minipage} \\
\midrule\noalign{}
\endhead
\bottomrule\noalign{}
\endlastfoot
\texttt{name} & Nome da configuração --- apenas identificação \\
\texttt{version} & Versão da sua configuração \\
\texttt{schema} & Versão do esquema do Continue --- manter
\texttt{v1} \\
\texttt{models} & Lista de modelos disponíveis para o chat \\
\texttt{name} (dentro de models) & Nome exibido no dropdown do
Continue \\
\texttt{provider} & Quem fornece o modelo (\texttt{ollama} para uso
local) \\
\texttt{model} & Nome exato do modelo conforme aparece no
\texttt{ollama\ list} \\
\texttt{apiBase} & Endereço da API do Ollama (\texttt{localhost:11434}
para uso local) \\
\texttt{allowAnonymousTelemetry} & Envio de dados anônimos de uso ---
recomenda-se \texttt{false} \\
\end{longtable}

\subsection{Principais atalhos do Continue no VS
Code}\label{principais-atalhos-do-continue-no-vs-code}

\begin{longtable}[]{@{}ll@{}}
\toprule\noalign{}
Atalho & Ação \\
\midrule\noalign{}
\endhead
\bottomrule\noalign{}
\endlastfoot
\texttt{Ctrl+Alt+I} & Abre o painel de chat \\
\texttt{Ctrl+Shift+L} & Envia o código selecionado para o chat \\
\texttt{Ctrl+I} & Edita o código selecionado inline com IA \\
\texttt{Tab} & Aceita sugestão de autocomplete \\
\texttt{Esc} & Rejeita sugestão de autocomplete \\
\end{longtable}

\subsection{Referenciando contexto no
chat}\label{referenciando-contexto-no-chat}

O Continue permite referenciar partes do projeto diretamente no chat:

\begin{longtable}[]{@{}ll@{}}
\toprule\noalign{}
Comando & O que referencia \\
\midrule\noalign{}
\endhead
\bottomrule\noalign{}
\endlastfoot
\texttt{@arquivo.py} & Um arquivo específico do projeto \\
\texttt{@codebase} & Todo o projeto \\
\texttt{@terminal} & A saída do terminal integrado \\
\texttt{@problems} & Os erros e avisos do VS Code \\
\end{longtable}

\begin{center}\rule{0.5\linewidth}{0.5pt}\end{center}

\section{Criando um Servidor de IA com
Ollama}\label{criando-um-servidor-de-ia-com-ollama}

\subsection{Conceito}\label{conceito}

O Ollama expõe uma API REST local na porta \texttt{11434}. Por padrão,
essa API só aceita conexões do próprio computador (\texttt{localhost}).
É possível configurá-lo para aceitar conexões de outros dispositivos na
rede local, transformando seu computador --- ou um dispositivo embarcado
como uma Jetson Nano --- em um servidor de IA compartilhado.

\begin{verbatim}
Servidor (Ollama)                    Clientes na rede
┌─────────────────────┐              ┌──────────────────┐
│ Modelo rodando      │ ──────────►  │ Notebook         │
│ API: 0.0.0.0:11434  │   rede local │ Desktop          │
│ IP: 192.168.1.100   │              │ Celular          │
└─────────────────────┘              └──────────────────┘
\end{verbatim}

\subsection{Configuração no Linux
(Ubuntu)}\label{configurauxe7uxe3o-no-linux-ubuntu}

\textbf{1. Editar o serviço do Ollama}

\begin{Shaded}
\begin{Highlighting}[]
\FunctionTok{sudo}\NormalTok{ systemctl edit ollama.service}
\end{Highlighting}
\end{Shaded}

\textbf{2. Adicionar a variável de ambiente}

\begin{Shaded}
\begin{Highlighting}[]
\ExtensionTok{[Service]}
\VariableTok{Environment}\OperatorTok{=}\StringTok{"OLLAMA\_HOST=0.0.0.0"}
\end{Highlighting}
\end{Shaded}

\textbf{3. Reiniciar o serviço}

\begin{Shaded}
\begin{Highlighting}[]
\FunctionTok{sudo}\NormalTok{ systemctl daemon{-}reload}
\FunctionTok{sudo}\NormalTok{ systemctl restart ollama}
\end{Highlighting}
\end{Shaded}

\textbf{4. Verificar se está aceitando conexões externas}

\begin{Shaded}
\begin{Highlighting}[]
\CommentTok{\# Substitua pelo IP do servidor na sua rede}
\ExtensionTok{curl}\NormalTok{ http://192.168.1.100:11434}
\end{Highlighting}
\end{Shaded}

\subsection{Configuração no
Windows}\label{configurauxe7uxe3o-no-windows}

No PowerShell como administrador:

\begin{Shaded}
\begin{Highlighting}[]
\CommentTok{\# Definir variável de ambiente permanente}
\ExtensionTok{[System.Environment]::SetEnvironmentVariable}\ErrorTok{(}
  \StringTok{"OLLAMA\_HOST"}\ExtensionTok{,}
  \StringTok{"0.0.0.0"}\ExtensionTok{,}
  \StringTok{"Machine"}
\KeywordTok{)}

\CommentTok{\# Reiniciar o serviço do Ollama}
\ExtensionTok{Stop{-}Service}\NormalTok{ ollama}
\ExtensionTok{Start{-}Service}\NormalTok{ ollama}
\end{Highlighting}
\end{Shaded}

\subsection{Acessando o servidor
remotamente}\label{acessando-o-servidor-remotamente}

Após configurar o servidor, os outros dispositivos da rede acessam o
modelo apontando para o IP do servidor:

\subsubsection{Via Python}\label{via-python}

\begin{Shaded}
\begin{Highlighting}[]
\ImportTok{import}\NormalTok{ ollama}

\CommentTok{\# Apontando para o servidor remoto}
\NormalTok{client }\OperatorTok{=}\NormalTok{ ollama.Client(host}\OperatorTok{=}\StringTok{\textquotesingle{}http://192.168.1.100:11434\textquotesingle{}}\NormalTok{)}

\NormalTok{response }\OperatorTok{=}\NormalTok{ client.chat(}
\NormalTok{    model}\OperatorTok{=}\StringTok{\textquotesingle{}qwen2.5{-}coder:0.5b\textquotesingle{}}\NormalTok{,}
\NormalTok{    messages}\OperatorTok{=}\NormalTok{[\{}\StringTok{\textquotesingle{}role\textquotesingle{}}\NormalTok{: }\StringTok{\textquotesingle{}user\textquotesingle{}}\NormalTok{, }\StringTok{\textquotesingle{}content\textquotesingle{}}\NormalTok{: }\StringTok{\textquotesingle{}Olá! Você está funcionando?\textquotesingle{}}\NormalTok{\}]}
\NormalTok{)}

\BuiltInTok{print}\NormalTok{(response[}\StringTok{\textquotesingle{}message\textquotesingle{}}\NormalTok{][}\StringTok{\textquotesingle{}content\textquotesingle{}}\NormalTok{])}
\end{Highlighting}
\end{Shaded}

\subsubsection{Via Continue (VS Code em outro
computador)}\label{via-continue-vs-code-em-outro-computador}

Basta apontar o \texttt{apiBase} no \texttt{config.yaml} para o IP do
servidor:

\begin{Shaded}
\begin{Highlighting}[]
\FunctionTok{models}\KeywordTok{:}
\AttributeTok{  }\KeywordTok{{-}}\AttributeTok{ }\FunctionTok{name}\KeywordTok{:}\AttributeTok{ Qwen no Servidor}
\AttributeTok{    }\FunctionTok{provider}\KeywordTok{:}\AttributeTok{ ollama}
\AttributeTok{    }\FunctionTok{model}\KeywordTok{:}\AttributeTok{ qwen2.5{-}coder:0.5b}
\AttributeTok{    }\FunctionTok{apiBase}\KeywordTok{:}\AttributeTok{ http://192.168.1.100:11434}
\end{Highlighting}
\end{Shaded}

\subsubsection{Via API REST (qualquer
dispositivo)}\label{via-api-rest-qualquer-dispositivo}

\begin{Shaded}
\begin{Highlighting}[]
\ExtensionTok{curl}\NormalTok{ http://192.168.1.100:11434/api/chat }\DataTypeTok{\textbackslash{}}
  \AttributeTok{{-}H} \StringTok{"Content{-}Type: application/json"} \DataTypeTok{\textbackslash{}}
  \AttributeTok{{-}d} \StringTok{\textquotesingle{}\{}
\StringTok{    "model": "qwen2.5{-}coder:0.5b",}
\StringTok{    "messages": [\{"role": "user", "content": "Olá!"\}],}
\StringTok{    "stream": false}
\StringTok{  \}\textquotesingle{}}
\end{Highlighting}
\end{Shaded}

\subsection{Servidor em dispositivo embarcado: Jetson
Nano}\label{servidor-em-dispositivo-embarcado-jetson-nano}

A NVIDIA Jetson Nano é uma opção interessante para servidor de IA de
baixo consumo energético (5--10W), porém o Ollama oficial não suporta
sua arquitetura ARM com GPU Maxwell nativamente. A alternativa
recomendada é o \textbf{llama.cpp}:

\begin{Shaded}
\begin{Highlighting}[]
\CommentTok{\# Clonar e compilar o llama.cpp com suporte CUDA}
\FunctionTok{git}\NormalTok{ clone https://github.com/ggerganov/llama.cpp}
\BuiltInTok{cd}\NormalTok{ llama.cpp}
\FunctionTok{make}\NormalTok{ LLAMA\_CUBLAS=1}

\CommentTok{\# Iniciar o servidor}
\ExtensionTok{./server} \DataTypeTok{\textbackslash{}}
  \AttributeTok{{-}m}\NormalTok{ models/tinyllama.gguf }\DataTypeTok{\textbackslash{}}
  \AttributeTok{{-}{-}host}\NormalTok{ 0.0.0.0 }\DataTypeTok{\textbackslash{}}
  \AttributeTok{{-}{-}port}\NormalTok{ 11434}
\end{Highlighting}
\end{Shaded}

Modelos recomendados para a Jetson Nano (4 GB RAM compartilhada):

\begin{longtable}[]{@{}ll@{}}
\toprule\noalign{}
Modelo & Viabilidade \\
\midrule\noalign{}
\endhead
\bottomrule\noalign{}
\endlastfoot
\texttt{tinyllama} & ✅ Roda bem \\
\texttt{qwen2.5-coder:0.5b} & ✅ Roda bem \\
\texttt{llama3.2:1b} & ⚠️ Aceitável \\
\texttt{phi3:mini} & ⚠️ Lento \\
\texttt{llama3.2:3b} & ❌ Muito pesado \\
\end{longtable}

\subsection{Boas práticas de segurança para o
servidor}\label{boas-pruxe1ticas-de-seguranuxe7a-para-o-servidor}

\begin{itemize}
\tightlist
\item
  \textbf{Use apenas em rede local confiável} --- não exponha a porta
  11434 para a internet sem autenticação
\item
  \textbf{Configure um firewall} para limitar o acesso apenas aos
  dispositivos autorizados
\item
  \textbf{Monitore o consumo de recursos} --- múltiplos clientes
  simultâneos podem sobrecarregar o servidor
\item
  Se precisar de acesso remoto seguro, use uma \textbf{VPN} ou
  \textbf{túnel SSH}
\end{itemize}

\begin{center}\rule{0.5\linewidth}{0.5pt}\end{center}

\section{Referências e Recursos}\label{referuxeancias-e-recursos}

\begin{itemize}
\tightlist
\item
  \href{https://ollama.com}{Ollama --- Site oficial}
\item
  \href{https://github.com/ollama/ollama}{Ollama --- Repositório no
  GitHub}
\item
  \href{https://continue.dev}{Continue --- Site oficial}
\item
  \href{https://github.com/ggerganov/llama.cpp}{llama.cpp ---
  Repositório no GitHub}
\item
  \href{https://github.com/ollama/ollama-python}{Biblioteca Python do
  Ollama}
\item
  \href{https://ollama.com/library}{Modelos disponíveis no Ollama}
\end{itemize}

\begin{center}\rule{0.5\linewidth}{0.5pt}\end{center}

\emph{Tutorial elaborado com base em experiência prática com Ubuntu
22.04 LTS e Windows 11.}

\end{document}
